\documentclass[10pt,a4paper]{amsart}
\usepackage[margin=19mm]{geometry}
\usepackage{amsmath,amssymb,mathtools}
\usepackage[scr=rsfs,cal=euler]{mathalfa}
\usepackage{xfrac}
\usepackage[unicode,hidelinks,bookmarks=false]{hyperref}
\newcommand{\ie}{i.e.\ }
\newcommand{\cf}{cf.\ }
\newcommand{\resp}{respectively\ }
\newcommand{\Subsection}{\subsection}
\providecommand{\nobreakdash}{\nobreak\hbox{-}}
\setlength{\emergencystretch}{2em}
\multlinegap0pt
\usepackage{amssymb}
\usepackage[shortlabels]{enumitem}
\usepackage{tikz}
\newtheorem{theorem}{Theorem}
\newtheorem{cor}{Corollary}
\title{\vspace{-1.2cm} \bf Local Rigidity of the Bergman Metric and\\ of the K\"ahler Carath\'eodory Metric}
\author{Robert Xin Dong, Ruoyi Wang, Bun Wong}
\date{Source: arXiv:2408.09572v3 (CC BY 4.0)}
\begin{document}
\maketitle

\noindent\emph{Source and licence.} \vspace{-1.2cm} \bf Local Rigidity of the Bergman Metric and\\ of the K\"ahler Carath\'eodory Metric. Version arXiv:2408.09572v3 (CC BY 4.0), distributed by arXiv under the Creative Commons Attribution 4.0 International licence, http://creativecommons.org/licenses/by/4.0/, as recorded in the arXiv licence metadata for this exact version and shown on the abstract page https://arxiv.org/abs/2408.09572v3. Full licence notice: \texttt{licenses/arxiv-2408.09572-CC-BY-4.0.txt}.\par\smallskip

\noindent\textit{Editorial scope.} Counted unit: the article's cor cor (source0.tex lines 319--331). The article's complete original proof of it is delivered with it (source0.tex lines 937--957, one \texttt{proof} environment of 21 lines, which the article places away from the statement). No mathematics is written, repaired or re-derived: the statement and the proof are the article's own lines, in the article's own notation, and no retained line is substituted or re-flowed.  Retained uncounted as the source-local context the proof needs: lines 232--234; lines 267--298.  Nothing was omitted from any retained span, so the package carries no omission record.  No retained line contains a citation, so the wrapper carries no bibliography and no bibliography entry.  No retained line contains a figure environment, an image-inclusion command or drawing code, so no image is delivered and the package declares no asset dependency.  Editorial formatting only, in addition: an AMS \texttt{amsart} preamble that restates the article's own declarations and macros used by the retained lines, namely the environments \texttt{theorem}, \texttt{cor} on separate counters, together with the standard packages those lines need.  The retained environments print in this document as Theorem 1, Corollary 1 in order of appearance, whereas the article prints the same items with the numbering of its own full document.  The complete native source of the article as distributed in the arXiv e-print for this exact version is bundled unchanged as \texttt{sources/163930/source0.tex}.
\begin{equation} \label{sharp}
L(\Omega) B_\Omega \geq C_\Omega.
\end{equation}

\begin{theorem} \label{rigidity}


 

\item [{$\bf A$.}] Let $X$ be a hyperbolic Riemann surface.
Assume that the Bergman metric $B(z)|dz|^2$ is complete with holomorphic sectional curvature no less than \(-1\), 
and 
the Carath\'eodory metric is written as $c^2_B(z) |dz|^2$,
where $z$ is a local coordinate.
Then, it holds on $X$ that
 $$
{B} \geq   2c^2_B.
 $$
Moreover, $X$ is biholomorphic to a disc if and only if there exists some $z_0 \in X$ such that
$$
B(z_0) =   2 c^2_B(z_0).
$$



 \item [{$\bf B$.}]

Let $\Omega \subset \mathbb{C}^n$ be a  domain whose
  universal cover  is biholomorphic to a Euclidean ball. Assume that  the   Bergman metric $B^2_{\Omega}$  is complete on $\Omega$ with holomorphic sectional curvature no less than $-2(L(\Omega))^2$. Then, $\Omega$ is biholomorphic to a ball  if and only if   there exists some  $z_0 \in \Omega$ such that 
$$
 L(\Omega) B_\Omega (z_0, v) = C_{\Omega} (z_0, v), \quad \text{for all\,} \, v\in T_{z_0}(\Omega). 
$$
   
 

\end{theorem}  

\begin{cor} \label{cor}




{\it Let $\Omega$ be a bounded strictly pseudoconvex domain  with spherical boundary in $\mathbb{C}^n$, $n \geq 2$. Then 
$\Omega$ is not biholomorphic to a Euclidean ball if and only if either the 
holomorphic sectional curvature of the Bergman metric   in some direction is strictly less than 
$-\frac{2}{n+1}$ or the sharp Lu's inequality \eqref{sharp} is strict at all points and for some directions.}



\end{cor}  

  \begin{proof} [{\bf Proof of Corollary \ref{cor}}]
  
We will prove a negated version of the corollary, namely,
$\Omega$ is biholomorphic to a Euclidean ball if and only if the holomorphic sectional curvature of the Bergman metric is no less than  $-\frac{2}{n+1}$ and equality is attained in the sharp Lu's inequality at a point in $\Omega$.

 \medskip



The necessity holds true because for a domain 
biholomorphic to a Euclidean ball in $\mathbb{C}^n$, the 
holomorphic sectional curvature of the Bergman metric is identically equal to  $-\frac{2}{n+1}$
and  the sharp Lu's inequality \eqref{sharp} becomes equality everywhere.


\medskip

For the sufficiency, if the holomorphic sectional curvature of the Bergman metric is no less than  $-\frac{2}{n+1}$, then  \(L(\Omega) = \frac{1}{\sqrt{n+1}}\) following from the Ahlfors-Schwarz lemma. So, the holomorphic sectional curvature is no less than $-2(L(\Omega))^2$. If equality is attained in the sharp Lu's inequality at a point in $\Omega$, then  we conclude by Theorem \ref{rigidity}, Part {$\bf B$}, that $\Omega$ is biholomorphic to a  ball.


\end{proof}

\end{document}
